% ============================================================
% Cheat sheet, argomento 12: il diagramma di attività (dispensa, capitolo 9)
% ============================================================
\argomento{Il diagramma di attività}{Dispensa: capitolo 9. Un flusso di lavoro o un algoritmo, anche
concorrente}[Diagramma di attività]

\voce{azione}{Azione, inizio e fine}{rettangoli arrotondati collegati da frecce}
  {\begin{tikzpicture}
     \node[iniziale] (i) at (0,0) {};
     \node[azione, minimum width=1.9cm] (a) at (1.8,0) {Ricevi ordine};
     \node[azione, minimum width=1.6cm] (b) at (4.1,0) {Spedisci};
     \node[finale] (f) at (5.6,0) {};
     \draw[tr] (i) -- (a); \draw[tr] (a) -- (b); \draw[tr] (b) -- (f);
   \end{tikzpicture}}
  {Un'\textbf{azione} è un passo di elaborazione (un verbo); le frecce sono il \textbf{flusso}: si entra quando
   l'azione inizia, si esce quando finisce. Il nodo finale chiude \textbf{l'intera} attività, anche gli altri flussi
   ancora in corso.}

\voce{finaleflusso}{Fine di un flusso}{cerchio con la croce \oltre}
  {\begin{tikzpicture}
     \node[azione, minimum width=1.9cm] (a) at (0,0) {Invia avviso};
     \node[circle, draw=black, fill=white, line width=0.8pt, inner sep=0pt, minimum size=0.36cm] (f) at (2.0,0) {};
     \draw[line width=0.7pt] ($(f.center)+(-0.12,-0.12)$) -- ($(f.center)+(0.12,0.12)$)
                             ($(f.center)+(-0.12,0.12)$) -- ($(f.center)+(0.12,-0.12)$);
     \draw[tr] (a) -- (f);
   \end{tikzpicture}}
  {Termina \textbf{solo quel flusso}, non tutta l'attività: utile dopo un fork, quando un ramo finisce prima degli
   altri senza doverli interrompere.}

\voce{branch}{Branch e merge}{il rombo: un solo ramo}
  {\begin{tikzpicture}
     \node[decisione] (d) at (0,0) {};
     \node[azione, minimum width=1.5cm] (a) at (-1.4,-1.0) {Espressa};
     \node[azione, minimum width=1.5cm] (b) at (1.4,-1.0) {Normale};
     \node[decisione] (m) at (0,-2.0) {};
     \draw[tr] (0,0.6) -- (d);
     \draw[tr] (d) -| node[etich, above, pos=0.25]{[prioritario]} (a);
     \draw[tr] (d) -| node[etich, above, pos=0.25]{[else]} (b);
     \draw[tr] (a) |- (m); \draw[tr] (b) |- (m);
     \draw[tr] (m) -- ++(0,-0.5);
   \end{tikzpicture}}
  {Il \textbf{branch} è un \code{if}: un arco entra, più archi escono, ognuno con una \textbf{guardia} fra quadre
   (mutuamente esclusive, \code{[else]} per il resto). Il \textbf{merge} è la fine dell'\code{if}: più archi entrano,
   ne esce uno. Un arco che torna indietro a un rombo è un \textbf{ciclo}.}

\voce{fork}{Fork e join}{la barra: tutti i rami, insieme}
  {\begin{tikzpicture}
     \node[barra, minimum width=3.0cm] (f) at (0,0) {};
     \node[azione, minimum width=1.5cm] (a) at (-1.0,-1.0) {Prepara};
     \node[azione, minimum width=1.5cm] (b) at (1.0,-1.0) {Fattura};
     \node[barra, minimum width=3.0cm] (j) at (0,-2.0) {};
     \draw[tr] (0,0.6) -- (f);
     \draw[tr] (-1.0,-0.05) -- (a); \draw[tr] (1.0,-0.05) -- (b);
     \draw[tr] (a) -- (-1.0,-1.95); \draw[tr] (b) -- (1.0,-1.95);
     \draw[tr] (j) -- ++(0,-0.55);
   \end{tikzpicture}}
  {Il \textbf{fork} avvia più flussi \textbf{concorrenti} (thread): dopo la barra si percorrono \textbf{tutti} i rami. Il
   \textbf{join} aspetta che arrivino \textbf{tutti} prima di proseguire. Il rombo sceglie, la barra sdoppia.}

\voce{swimlane}{Swimlane}{corsie: chi fa che cosa}
  {\begin{tikzpicture}
     \draw[line width=0.7pt] (0,0) rectangle (5.2,-1.8);
     \draw[line width=0.7pt] (2.6,0) -- (2.6,-1.8);
     \draw[line width=0.5pt] (0,-0.35) -- (5.2,-0.35);
     \node[font=\tiny\bfseries] at (1.3,-0.18) {Cliente};
     \node[font=\tiny\bfseries] at (3.9,-0.18) {Negozio};
     \node[azione, minimum width=1.7cm, minimum height=0.55cm] (o) at (1.3,-0.95) {Ordina};
     \node[azione, minimum width=1.7cm, minimum height=0.55cm] (s) at (3.9,-0.95) {Spedisce};
     \draw[tr] (o) -- (s);
   \end{tikzpicture}}
  {Il diagramma si divide in \textbf{corsie}, una per classe, ruolo, reparto o sistema; ogni azione sta nella corsia di
   chi la esegue. Così il diagramma di attività mostra anche la collaborazione fra oggetti diversi.}

\voce{nodooggetto}{Nodo oggetto e datastore}{l'oggetto che passa da un'azione all'altra}
  {\begin{tikzpicture}
     \node[azione, minimum width=1.6cm] (a) at (0,0) {Scrive};
     \node[draw=black!75, fill=white, line width=0.7pt, font=\scriptsize, align=center, minimum height=0.6cm] (o) at (2.2,0) {Compito\\{}[consegnato]};
     \node[azione, minimum width=1.6cm] (b) at (4.4,0) {Corregge};
     \node[draw=black!75, fill=white, line width=0.7pt, font=\scriptsize, align=center] (d) at (4.4,-1.1) {{\tiny«datastore»}\\Registro voti};
     \draw[tr] (a) -- (o); \draw[tr] (o) -- (b); \draw[tr] (b) -- (d);
   \end{tikzpicture}}
  {Un rettangolo con il nome di un oggetto e, fra quadre, il suo \textbf{stato} in quel punto del flusso. Il
   \textbf{«datastore»} è un archivio persistente: i dati restano anche quando l'attività finisce.}

\voce{macro}{Macro-attività}{il rastrello: c'è un sotto-diagramma}
  {\begin{tikzpicture}
     \node[azione, minimum width=2.4cm, minimum height=0.8cm, line width=1.2pt, text depth=6pt] (d) at (0,0) {Consegna ordine};
     \begin{scope}[shift={($(d.south east)+(-0.3,0.08)$)}, line width=0.6pt]
       \draw (0,0) -- (0,0.18); \draw (-0.1,0.09) -- (0.1,0.09); \draw (-0.1,0.09) -- (-0.1,0); \draw (0.1,0.09) -- (0.1,0);
     \end{scope}
     \node[draw=black!70, fill=white, font=\tiny, minimum width=0.9cm, minimum height=0.45cm] (p) at (3.1,0) {Ordine};
     \node[font=\tiny, anchor=north] at (3.1,-0.3) {parametro};
   \end{tikzpicture}}
  {Un compito complesso diventa una sola azione, con il simbolo del \textbf{rastrello}, e si dettaglia in un altro
   diagramma. Il dettaglio è un'attività con un nome e dei \textbf{parametri}: rettangoli sul bordo per i dati in
   ingresso e in uscita.}

\voce{segnali}{Segnali}{tempo, ricevuto, inviato}
  {\begin{tikzpicture}[sig/.style={line width=0.8pt, draw=ph4!80!black, fill=ph4!8}]
     \draw[sig] (-0.25,0.3) -- (0.25,0.3) -- (-0.25,-0.3) -- (0.25,-0.3) -- cycle;
     \node[font=\tiny] at (0,-0.5) {tempo};
     \draw[sig] (0.9,0.3) -- (2.5,0.3) -- (2.5,-0.3) -- (0.9,-0.3) -- (1.2,0) -- cycle;
     \node[font=\tiny] at (1.75,0) {ricevuto};
     \draw[sig] (3.1,0.3) -- (4.5,0.3) -- (4.8,0) -- (4.5,-0.3) -- (3.1,-0.3) -- cycle;
     \node[font=\tiny] at (3.9,0) {inviato};
   \end{tikzpicture}}
  {Eventi \textbf{esterni} e \textbf{asincroni}: la \textbf{clessidra} è un segnale di tempo («due ore prima del volo»,
   «dopo 48 ore»); il pentagono con la rientranza è un segnale \textbf{ricevuto} (accept, «arriva il taxi»); quello a
   punta un segnale \textbf{inviato} (send).}

\begin{caso}{codiceattivita}{Dal codice al diagramma di attività}
\begin{lstlisting}[style=javacs]
import java.util.concurrent.CompletableFuture;

class GestioneOrdini {
    void gestisci(Ordine o) {
        // fork: i due flussi partono insieme
        CompletableFuture<Void> spedizione =
            CompletableFuture.runAsync(() -> {
                if (o.prioritario()) espressa(o);
                else normale(o);
            });
        CompletableFuture<Void> pagamento =
            CompletableFuture.runAsync(() -> {
                fattura(o);
                incassa(o);
            });
        // join: aspetta che finiscano entrambi
        CompletableFuture
            .allOf(spedizione, pagamento).join();
        chiudi(o);
    }
    void espressa(Ordine o) { }
    void normale(Ordine o) { }
    void fattura(Ordine o) { }
    void incassa(Ordine o) { }
    void chiudi(Ordine o) { }
}
class Ordine {
    boolean prioritario() { return false; }
}
\end{lstlisting}
\tcblower
\begin{tikzpicture}
  \node[iniziale] (i) at (0,0) {};
  \node[barra, minimum width=3.4cm] (fk) at (0,-0.7) {};
  \node[decisione] (d) at (-1.3,-1.5) {};
  \node[azione, minimum width=1.4cm] (es) at (-2.1,-2.5) {Espressa};
  \node[azione, minimum width=1.4cm] (no) at (-0.5,-2.5) {Normale};
  \node[decisione] (m) at (-1.3,-3.5) {};
  \node[azione, minimum width=1.5cm] (fa) at (1.5,-1.9) {Fattura};
  \node[azione, minimum width=1.5cm] (in) at (1.5,-3.1) {Incassa};
  \node[barra, minimum width=3.4cm] (jn) at (0,-4.3) {};
  \node[azione, minimum width=1.6cm] (ch) at (0,-5.1) {Chiudi};
  \node[finale] (f) at (0,-5.9) {};
  \draw[tr] (i) -- (fk);
  \draw[tr] (-1.3,-0.75) -- (d);
  \draw[tr] (1.5,-0.75) -- (fa);
  \draw[tr] (d) -| node[etich, above, pos=0.2]{[prioritario]} (es);
  \draw[tr] (d) -| node[etich, above, pos=0.2]{[else]} (no);
  \draw[tr] (es) |- (m); \draw[tr] (no) |- (m);
  \draw[tr] (m) -- (-1.3,-4.25);
  \draw[tr] (fa) -- (in);
  \draw[tr] (in) -- (1.5,-4.25);
  \draw[tr] (jn) -- (ch); \draw[tr] (ch) -- (f);
\end{tikzpicture}
\regola{Un \textbf{if}/\code{else} diventa branch e merge (rombi); due compiti avviati \textbf{in parallelo} diventano
un fork, e l'attesa che finiscano entrambi (\code{allOf(\ldots).join()}) un join (barre). Un \code{while} è un rombo
con un arco che torna indietro.}
\end{caso}

\begin{attenzione}{Errori tipici}
\begin{itemize}[nosep, leftmargin=*]
  \item Rombo al posto della barra (o il contrario): dopo il \textbf{rombo} si prende \textbf{un} ramo, dopo la
        \textbf{barra} \textbf{tutti}; il merge lascia passare il primo flusso, il join li aspetta tutti.
  \item Guardie che non si escludono, o che non coprono tutti i casi (manca \code{[else]}).
  \item Un nodo finale nel ramo di un fork: chiude anche gli altri rami (se non è voluto, serve la fine di flusso).
  \item Un flusso che si ferma nel nulla: ogni percorso deve arrivare a un nodo finale.
\end{itemize}
\end{attenzione}
